In our problems, a decision $z\in\mathcal{Z}$ must be chosen from a
discrete feasible set to minimize a cost that depends on outcomes $Y$ that are
unknown at decision time. We model $Y$ as a random variable produced from a
conditional distribution given observed features $X=x$. 
For example, in a later top-K site selection tasks each $z$ represents an indicator vector that picks $K$ sites for intervention out of many, while vector $y$ has an observed score for each site, with smaller values indicating desirable sites given $x$.
%hospitalized patient monitoring application, $z$ may represent a subset of patients selected to receive additional clinical attention from among all feasible subsets $\mathcal{Z}$, while $y$ indicates survival or death in next 24 hours for all patients, given patient-specific attributes $x$. 
The decision to do $z$ when the outcome is $y$ incurs a scalar cost $c(z,y)$. 

The decision-making problem is, given an $x$, to solve for the value(s) of $z$ with lowest expected cost,\begin{align}
%% OLD VERSION:
% z^\star(x)\;\in\;\arg\min_{z\in\mathcal{Z}} \mathbb{E} \! \left[c(z,Y)\mid X=x\right].
% \\ \notag 
%% MCH PROPOSED:
z^\star(x,\theta)\;:=\;\arg\min_{z\in\mathcal{Z}} \mathbb{E}_{p_{\theta}(Y | X=x)} \!\left[c(z,Y) \right].
\label{eq:decisions}
\end{align}
Here, the parameter $\theta \in \mathbb{R}^P$ defines a conditional distribution over $Y$ given $X{=}x$.
At decision-making time (``test time''), only features $x$ and parameter $\theta$ are available; the realization of $Y$ is
unknown until later.

The training problem is to estimate $\theta$ given observed data $\{(x_i,y_i)\}_{i=1}^n$ assumed drawn i.i.d.\ from the joint distribution
of $(X,Y)$.
We could learn $\theta$ to  maximize the likelihood $p_{\theta}(Y|X)$ given observed data.
However, end-to-end methods seek an estimated $\theta$ that delivers good decisions via Eq.~\eqref{eq:decisions} when used on never-before-seen contexts $X{=}x'$ drawn from the same distribution.
To solve this, training can be framed as a cost minimization problem: $\theta^* \gets \arg\!\min_{\theta} \sum_{i=1}^n c(z^*(x_i,\theta), y_i)$.


\textbf{Linear cost assumption.}
A common assumption~\citep{huang2024decisionfocusedPG} is that cost is a linear inner product $c(z, y) = 
\langle z, y\rangle$. Here $z$ and $y$ are vectors of size $D$ and  decisions simplify due to linearity of expectation: $z^*(x,\theta) = \arg\!\min_{z \in \mathcal{Z}} 
\langle z, \hat{y}_{\theta}(x) \rangle$, where $\hat{y}_{\theta}(x)$ denotes the mean of $Y$ under $p_{\theta}(Y | X=x)$.
Even here, the decision problem is  piecewise constant as a function of $\hat{y}_{\theta}$. Its flatness yields all-zero gradients which break gradient-based learning of $\theta$.

\textbf{DPO method.} The differentiable perturbed optimizer~\citep{berthetLearningDifferentiablePerturbed2020}, addresses this issue by smoothing the decision
map via random perturbations of the predictions. For noise
$\xi$ from a standard Normal or Gumbel distribution with noise-level $\sigma>0$, the perturbed decision is defined as
\begin{align}
z_\sigma(x, \theta)
\;:=\;
%\textstyle
\mathbb{E}_{\xi}[
~
\arg\min_{z\in\mathcal{Z}} \langle z,\, \hat{y}_{\theta}(x) + \sigma \xi\rangle
~].
\end{align}
In practice, we can estimate this perturbed decision vector of size $D$ and its gradient via Monte Carlo estimates. We focus on iid Gaussian noise and assume $M$ samples, each one denoted $\xi^{(m)} \sim \mathcal{N}(0, I_D)$. Define the decision obtained with each noise sample as as $\hat{z}^{(m)} {=} \arg\min_{z\in\mathcal{Z}} 
~\langle z,\, \hat{y}_{\theta}(x) + \sigma \xi^{(m)} \rangle$, then we can estimate the perturbed decision $\hat{z}$ and the Jacobian matrix $J$ needed for gradients as
\begin{align}
\hat{z}(x, \theta)
&=
\textstyle 
\frac{1}{M} \sum_{m=1}^M
\hat{z}^{(m)}
\\ \notag 
\hat{J}(x, \theta) &= 
\nabla_{\theta} \hat{y}(x,\theta) \Big[ 
\textstyle
\frac{1}{\sigma M} 
\sum_{m=1}^M \textsc{outer}(\hat{z}^{(m)}, \xi^{(m)})
\Big] 
\end{align}
This Jacobian estimator 
$\hat{J} \approx \nabla_{\theta} {z}_{\sigma}(x, \theta)$, a $P\times D$ matrix, comes directly from the chain rule and Prop. 3.1 of \citeauthor{berthetLearningDifferentiablePerturbed2020}~Noise distributions other than Gaussian require changes to the outer product term.

\textbf{DPO hyperparameters.}
DPO's most sensitive hyperparameter is noise-level $\sigma$, which controls a bias--smoothness tradeoff. Smaller $\sigma$ values yield high-variance or vanishing
gradients. We recover noise-free decisions $z^*(x,\theta)$ as $\sigma \rightarrow 0^+$. Larger $\sigma$ induce smoother decision maps at the cost of
biasing decisions away from $z^\star$. 
%Desireably, even with a single sample and small $\sigma$ the Jacobian will be non-zero.

As typical in Monte Carlo methods, DPO's estimators become more accurate as the number of samples $M$ increases: larger $M$ produce lower-variance estimates at the cost of an increased number of calls to the (often expensive) optimization solver that produces $\hat{z}^{(m)}$.

% at added computational cost. The number of samples $M$ controls the variance of the Monte Carlo approximation; a larger number of samples provides a lower-variance decision at the cost of an increased number of calls to the (often expensive) optimization solver.

% For either the ideal expectation or the M-sample estimator, as long as $\sigma$ is large enough that some different $z$ are chosen in the inner $\arg\!\min$, the decision will not suffer from the ``flatness'' problem, thus allowing informative (non-zero) gradients to be backpropagated through the entire computational graph.




%\subsection{Gradient SNR}

%Max Welling 2020 MRI NeurIPS, 2008 weight perturbation NeurIPS.